% formal/billiard-capacity-algorithm.tex — Lagrangian-product reduction
% Sources: thesis/lagrangian-product-algorithm.tex,
%   thesis/lagrangian-product-algorithm-proof.tex

\begin{unverified}
\begin{lemma}[Facet structure of Lagrangian products]\label{lem:lagrangian-facets}
% Source: thesis/lagrangian-product-algorithm-proof.tex
Let $K = K_q \times K_p$ be a Lagrangian product.
Let $K_q = \{x \in L_q : \langle a_k^q, x \rangle \leq 1\}$
with dual vertices
$a_1^q, \ldots, a_{F_q}^q \in L_q$,
and similarly $K_p$ with dual vertices
$a_1^p, \ldots, a_{F_p}^p \in L_p$.
Then $K$ has $F = F_q + F_p$ facets, partitioned into:
\begin{itemize}
\item \emph{$q$-facets} with dual vertices $a_i = (a_i^q, 0)$,
  for $i = 1, \ldots, F_q$;
\item \emph{$p$-facets} with dual vertices $a_j = (0, a_j^p)$,
  for $j = 1, \ldots, F_p$.
\end{itemize}
Every dual vertex lies entirely in $L_q$ or entirely in $L_p$.
\end{lemma}

\begin{proof}
A facet of $K_q \times K_p$ arises from a facet (edge) of
one factor and the full other factor: either
$\text{edge}_i(K_q) \times K_p$ or $K_q \times \text{edge}_j(K_p)$.
The defining constraint $\langle a_i^q, x_q \rangle \leq 1$
for edge~$i$ of~$K_q$ lifts to
$\langle (a_i^q, 0),\, x \rangle \leq 1$,
so the dual vertex is $(a_i^q, 0)$;
similarly for $p$-facets.
\end{proof}
\end{unverified}


\begin{unverified}
\begin{lemma}[Sigma structure for Lagrangian products]\label{lem:sigma-structure}
% Source: thesis/lagrangian-product-algorithm-proof.tex
Let $K = K_q \times K_p$ be a Lagrangian product, and let
$(\sigma, \tau)$ be a simple Reeb orbit on $\partial K$
with all $\tau_k > 0$.
Write the cyclic sequence $\sigma$ as a word in the alphabet
$\{Q, P\}$ by replacing each $q$-facet index with $Q$ and each
$p$-facet index with $P$.
Then:
\begin{enumerate}
\item[(i)] The word has the form
  $\bigl([Q \mid QQ]\,[P \mid PP]\bigr)^k$ for some $k \geq 1$:
  alternating blocks of $q$- and $p$-facets, each block of length
  one or two.

\item[(ii)] Within each block of length two, the two facets are
  adjacent in $G_K$
  (consecutive edges of $K_q$ or $K_p$ sharing a vertex).
\end{enumerate}
\end{lemma}

\begin{proof}
\textbf{Claim: no three consecutive same-type facets.}\quad
Suppose for contradiction that $\sigma$ contains three
consecutive $q$-facets $q_1, q_2, q_3$ (all with $\tau > 0$).
The Reeb vectors $R_{q_1}, R_{q_2}, R_{q_3}$ all lie in~$L_p$,
so $q$ is constant throughout the three segments.
The orbit at the breakpoint between $q_1$ and $q_2$ lies on
both facets, meaning $\langle a_{q_1}^q, q \rangle = 1$
and $\langle a_{q_2}^q, q \rangle = 1$: the point $q$
lies on both edges $1$ and $2$ of~$K_q$, hence at their
common vertex $v_{12}$.
Similarly, the breakpoint between $q_2$ and $q_3$ forces $q$
to lie at vertex $v_{23}$ (intersection of edges $2$ and $3$).
Since $q$ is constant, $v_{12} = v_{23}$.
But for a nondegenerate convex polygon, $v_{12}$ and $v_{23}$
are the two distinct endpoints of edge~$2$: contradiction.

The same argument applies to $p$-facets by symmetry.

\medskip\noindent
\textbf{Alternation and block structure.}\quad
Since no three consecutive facets have the same type, each maximal
run of same-type facets has length one or two.
% [TODO: JÖRN - alternation argument is imprecise: two adjacent length-1
%   q-runs would give only 2 consecutive q-facets, not 3. The correct
%   reason is that adjacent same-type maximal runs would merge.]
Consecutive runs must alternate between $q$ and $p$.
In a cyclic sequence, alternation forces equal block counts:
$k$ blocks of $q$-facets and $k$ blocks of $p$-facets.

\medskip\noindent
\textbf{Within-block adjacency.}\quad
Follows directly from the adjacency constraint:
consecutive facets in $\sigma$ share a breakpoint, hence are
adjacent in~$G_K$.
Two same-type adjacent facets correspond to edges sharing a
vertex in the 2D polygon.
\end{proof}
\end{unverified}


\begin{unverified}
\begin{theorem}[Billiard characterization of EHZ capacity]%
\label{thm:billiard-characterization}
% Source: thesis/lagrangian-product-algorithm-proof.tex
% Attribution: Artstein-Avidan \& Ostrover (2014); Rudolf (2022).
For a Lagrangian product $K = K_q \times K_p$:
\[
  c_{\mathrm{EHZ}}(K_q \times K_p)
  = \min\bigl\{ \ell_{K_p^\circ}(\eta) :
  \eta \text{ is a closed $K_p$-billiard trajectory in } K_q \bigr\},
\]
where $\ell_{K_p^\circ}(\eta) = \sum_i h_{K_p}(\Delta_i)$ is the
$K_p^\circ$-length and $\Delta_i = x_{i+1} - x_i$ are the edge vectors.
\end{theorem}
% Proof: Artstein-Avidan \& Ostrover (2014), Theorem~2.13; Rudolf (2022), Theorem~1.
% See also HKO2024, Theorem~2.2.
\end{unverified}


\begin{unverified}
\begin{theorem}[Bounce bound]%
\label{thm:bounce-bound}
% Source: thesis/lagrangian-product-algorithm-proof.tex
% Attribution: Bezdek \& Bezdek (2009); Rudolf (2022).
For convex polygons $K_q, K_p \subset \mathbb{R}^2$,
the minimum in Theorem~\ref{thm:billiard-characterization}
is achieved by a billiard trajectory with at most $3$ bounce points.
\end{theorem}
% Proof: Bezdek \& Bezdek (2009), Theorem~1.1 and Lemma~2.4; Rudolf (2022), Theorem~1.
\end{unverified}


\begin{unverified}
\begin{algorithm}[Billiard capacity]\label{alg:billiard}
% Source: thesis/lagrangian-product-algorithm.tex
Computes $c_{\mathrm{EHZ}}(K_q \times K_p)$
via Theorem~\ref{thm:billiard-characterization}.

\emph{Input:} Lagrangian product $K = K_q \times K_p$ in
$H$-representation.
\quad
\emph{Output:} $c_{\mathrm{EHZ}}(K)$.

\begin{enumerate}
\item \textbf{Classify} facets into $q$-type and $p$-type
  (Lemma~\ref{lem:lagrangian-facets}).
\item \textbf{Build blocks.}
  For each polygon ($K_q$ and $K_p$), enumerate all
  \emph{single-facet blocks} (one edge) and
  \emph{pair blocks} (two adjacent edges, both orderings).
% [TODO: JÖRN - k=1 is excluded here. The thesis proves this via
%   lem:k1-infeasible (closure condition forces antiparallel normals
%   on adjacent edges, impossible for convex polygons). Add citation or
%   copy the lemma.]
\item \textbf{Enumerate.}
  For $k = 2, 3$: for each selection of $k$ non-overlapping
  $q$-blocks and $k$ non-overlapping $p$-blocks, and each
  interleaving into the pattern
  $\bigl([Q \mid QQ]\,[P \mid PP]\bigr)^k$,
  flatten the blocks into a facet sequence~$\sigma$.
\item \textbf{Solve and filter.}
  For each $\sigma$: solve the KKT system,
  filter by $\beta_i > 0$ and $Q(\beta) > 0$, compute
  $A(S, \sigma) = \tfrac{1}{2}\, Q(\beta)^{-1}$.
\item \textbf{Return} $\min A(S, \sigma)$.
\end{enumerate}
\end{algorithm}
\end{unverified}
